\section{Motivation}
\label{sec:motivation}

\subsection{Why a logic learner at all}

Deep networks answer; they rarely account for their answers. In clinical decision support,
regulated industry and on-device inference, the account matters as much as the answer: a
cardiologist wants to know which measurement moved the prediction, an auditor wants a
decision rule they can read, and a microcontroller wants an inference path that is a handful
of bit operations rather than a matrix cascade.

The Tsetlin machine~\ttcite{1} occupies exactly that niche. It learns conjunctive clauses
over Boolean literals; clauses vote for or against each class; the class with the largest
vote sum wins (Figure~\ref{fig:tm-overview}). Three properties follow:

\begin{itemize}[leftmargin=1.25em, itemsep=0.25em, topsep=0.35em]
  \item \textbf{Legibility by construction.} A clause is a readable conjunction. There is no
    surrogate model, no saliency heuristic, and no gap between what is explained and what is
    executed. Global and local importance have closed-form expressions~\ttcite{4}.
  \item \textbf{Bit-level economy.} Inputs are Booleans, the learned state is small integers,
    and inference is an \textsc{and} over included literals plus a vote count --- which is
    why a substantial hardware literature has grown around it, from FPGA and ASIC inference
    engines to in-memory computing~\ttcite{15,16,17}.
  \item \textbf{Competitive accuracy} on the tasks where those properties are worth paying
    for: tabular and clinical data, keyword spotting, and small-image classification.
\end{itemize}

\begin{figure}[t]
\centering
\resizebox{\linewidth}{!}{% Figure: what a Tsetlin machine model is — literals, clauses, polarised votes, argmax.
\begin{tikzpicture}[font=\sffamily\small]

  % ---------------------------------------------------------------- input bits
  \node[tttitle, anchor=south west] at (-0.15,1.55) {Boolean input $x$};
  \foreach \i/\v in {0/1,1/0,2/1,3/1}{
    \pgfmathsetmacro{\yy}{1.15-\i*0.52}
    \ifnum\v=1
      \node[ttcell on] (x\i) at (0.25,\yy) {1};
    \else
      \node[ttcell off] (x\i) at (0.25,\yy) {0};
    \fi
    \node[ttlabel, anchor=east] at (-0.05,\yy) {$x_\i$};
  }

  % ---------------------------------------------------------------- literals
  \node[tttitle, anchor=south west] at (1.55,1.55) {Literals $[x,\lnot x]$};
  \foreach \i/\v in {0/1,1/0,2/1,3/1}{
    \pgfmathsetmacro{\yy}{1.15-\i*0.52}
    \ifnum\v=1 \node[ttcell on] (l\i) at (2.0,\yy) {1};
    \else \node[ttcell off] (l\i) at (2.0,\yy) {0}; \fi
  }
  \foreach \i/\v in {0/0,1/1,2/0,3/0}{
    \pgfmathsetmacro{\yy}{1.15-\i*0.52}
    \ifnum\v=1 \node[ttcell on] (n\i) at (2.52,\yy) {1};
    \else \node[ttcell off] (n\i) at (2.52,\yy) {0}; \fi
  }
  \node[ttlabel] at (2.0,-1.28) {$x_k$};
  \node[ttlabel] at (2.56,-1.28) {$\lnot x_k$};

  \draw[ttarrow] (0.62,-0.05) -- (1.72,-0.05);

  % ---------------------------------------------------------------- clauses
  \node[tttitle, anchor=south west] at (3.5,1.55) {Clauses (AND of included literals)};
  \node[ttnode, anchor=west, text width=4.0cm] (c1) at (3.5,1.15)
    {$C_1: x_0 \land \lnot x_1$};
  \node[ttnode, anchor=west, text width=4.0cm] (c2) at (3.5,0.38)
    {$C_2: x_2 \land x_3$};
  \node[ttnode plain, anchor=west, text width=4.0cm] (c3) at (3.5,-0.42)
    {$C_3: x_1 \land x_3$};
  \node[ttnode, anchor=west, text width=4.0cm] (c4) at (3.5,-1.18)
    {$C_4: \lnot x_0 \land x_2$};

  \node[ttlabel, anchor=west, text=ttorange] at (7.62,1.15) {matches};
  \node[ttlabel, anchor=west, text=ttorange] at (7.62,0.38) {matches};
  \node[ttlabel, anchor=west, text=ttgray] at (7.62,-0.42) {no match};
  \node[ttlabel, anchor=west, text=ttorange] at (7.62,-1.18) {matches};

  \draw[ttarrow] (2.88,-0.05) -- (3.36,-0.05);

  % ---------------------------------------------------------------- votes
  \node[tttitle, anchor=south] at (10.55,1.55) {Polarised votes};
  \node[anchor=west, font=\small, text=ttdeep] at (9.55,1.15) {$+1$};
  \node[anchor=west, font=\small, text=ttdeep] at (9.55,0.38) {$+1$};
  \node[anchor=west, font=\small, text=ttgray] at (9.55,-0.42) {$-1$};
  \node[anchor=west, font=\small, text=ttdeep] at (9.55,-1.18) {$-1$};

  \node[ttnode dark, anchor=west, text width=2.55cm] (v0) at (10.25,0.72)
    {class 0: $v_0 = +2$};
  \node[ttnode, anchor=west, text width=2.55cm] (v1) at (10.25,-0.75)
    {class 1: $v_1 = -1$};

  \draw[ttarrow] (9.95,1.15) .. controls (10.05,1.15) and (10.05,0.95) .. (10.21,0.9);
  \draw[ttarrow] (9.95,0.38) .. controls (10.05,0.38) and (10.05,0.6) .. (10.21,0.62);
  \draw[ttarrow, draw=ttgray] (9.95,-0.42) .. controls (10.05,-0.42) and (10.05,-0.6) .. (10.21,-0.62);
  \draw[ttarrow] (9.95,-1.18) .. controls (10.05,-1.18) and (10.05,-0.95) .. (10.21,-0.9);

  % ---------------------------------------------------------------- prediction
  \node[ttnode dark, anchor=west] (pred) at (13.35,-0.02)
    {$\hat y = \arg\max_k \mathrm{clip}(v_k,-T,T)$\\[1pt] $=\;$\textbf{class 0}};
  \draw[ttarrow] (12.92,0.72) -- (13.2,0.15);
  \draw[ttarrow] (12.92,-0.75) -- (13.2,-0.2);

  % ---------------------------------------------------------------- footer note
  \node[ttlabel, anchor=north west, text width=15.5cm, align=left] at (-0.15,-1.75)
    {Every clause is a conjunction the reader can check by hand; the prediction is a signed
     vote count over the clauses that matched. The model is its own explanation, and
     inference is an \textsc{and} plus an addition.};
\end{tikzpicture}
}
\caption{A Tsetlin machine classifier at inference time. Features are expanded into literals
  $[x,\lnot x]$; each clause is a conjunction of the literals its automata have decided to
  \emph{include}; matching clauses cast signed votes; the argmax over clipped vote sums is
  the prediction.}
\label{fig:tm-overview}
\end{figure}

\subsection{Three obstacles in practice}

The motivation for \ttlibname came from trying to use Tsetlin machines in a medical-imaging
research group, where the surrounding stack is PyTorch and the data are large. Three
obstacles showed up immediately, and none of them is about the algorithm.

\paragraph{1.~The learning rule looks hostile to a GPU.}
The classical algorithm is defined per example: evaluate the clauses, draw feedback,
\emph{update the automata}, then move to the next example. The update is stochastic and
touches every (clause, literal) counter, so a naive implementation performs one small,
serial, random-number-heavy pass per training example. On a device with thousands of cores
and a kernel-launch cost, that is the worst possible shape of work. A GPU implementation
therefore requires reformulating the update, not merely porting it.

\paragraph{2.~The tooling lives outside the ecosystem.}
Existing implementations are compiled kernels behind thin wrappers
(Section~\ref{sec:existing}). They are fast and faithful, but they sit beside a PyTorch
pipeline rather than inside it. Moving a model to a GPU, checkpointing it mid-training,
feeding it from a \code{DataLoader} that already augments on-device, or profiling it next to
the CNN it is being compared against --- each requires bespoke glue, and the data usually
makes a round trip through host memory and NumPy on the way.

\paragraph{3.~Everything around the model is re-invented per project.}
A Tsetlin machine consumes Booleans, so every project starts by writing a Booleanizer:
thermometer thresholds for tabular columns, adaptive thresholding for images, bit planes for
integers. Then a training loop, then early stopping, then a rule printer, then a clause
visualiser. This is the unglamorous 80\% of a Tsetlin machine project, it is rewritten
constantly, and small differences in it (which thresholds, which threshold strategy) move
accuracy more than the model hyper-parameters do.

\begin{ttcallout}[The design question]
Can the Tsetlin machine be expressed so that it is \emph{ordinary} inside PyTorch --- same
device semantics, same checkpoints, same data plumbing --- without giving up fidelity to the
published algorithm? \ttlibname is the answer to that question: yes, if the learning update
is reformulated from per-example state mutation into per-batch \emph{event counting}
(Section~\ref{sec:batched}), and if the exact per-example path is kept available for when
fidelity matters more than speed.
\end{ttcallout}
